\section{Chapter~\ref{sec:pain}. Pain}

Pain sensors, the two wires, the gate in the spinal cord, the descending volume knob, sensitization, and the capture of attention have been measured directly; the gate and sensitization formulas are simplified models of the book; the rule by which the machine chooses the volume of the alarm is a hypothesis.

{\footnotesize
\setlength{\tabcolsep}{3pt}\renewcommand{\arraystretch}{1.15}
\begin{longtable}{>{\raggedright}p{0.5cm}>{\raggedright}p{4.0cm}>{\raggedright}p{5.6cm}>{\raggedright}p{2.3cm}>{\raggedright\arraybackslash}p{1.7cm}}
\toprule
No. & Claim & Experiment and what was measured & Source & Status \\
\midrule\endfirsthead
\toprule
No. & Claim & Experiment and what was measured & Source & Status \\
\midrule\endhead
1 & High threshold: heat thresholds of different pain sensors range from $41^\circ$ to $46$--$48^\circ$ & Single-fiber recordings. Monkey skin: median threshold of the slow (C) sensors $41^\circ$, of fast type~II sensors $46^\circ$, of type~I above $53^\circ$. Human skin: C sensors that also respond to pressure $40.7^\circ$, those that do not respond to pressure $48^\circ$. & \cite{Treede1995heat, Weidner1999cfib} & measured \\
2 & Pain sensors adapt far less than touch sensors and become more sensitive after mild injury; a weakening after strong heating is measured in one type & A mild skin burn ($50^\circ$, $100\,\text{s}$): after $5$--$10$~min the pain threshold in humans and the threshold of C sensors in the monkey are lowered by $2$--$6^\circ$; other skin sensors show no such shift. In fast type~II sensors the response after strong heating is suppressed. The comparison of adaptation with touch sensors is a general estimate, not backed here by a separate experiment. & \cite{LaMotte1982hyper, Treede1995heat} & measured \\
3 & Two wires: fast $5$--$30$~m/s, slow $0.5$--$2$~m/s; arrival time $t=L/v$~\eqref{eq:paintime} & Conduction velocity: fast pain sensors of the monkey average $15$ and $25$~m/s (two types); human C sensors $0.8$--$1.0$~m/s. For $L=1$~m this gives tens of milliseconds and about a second. & \cite{Treede1995heat, Weidner1999cfib} & measured \\
4 & Pain arrives twice: first fast and precise, then dull and lasting & Reaction time to heating of the human forearm: from $1100$ to $400\ms$ as the temperature rises; strong heating of hairy skin gives double pain, of the palm does not. First pain can be carried only by fibers faster than $6$~m/s; by latency it corresponds to fast type~II sensors, which are absent from the palm. The division of roles (``where'' and ``how bad'') is an interpretation. & \cite{CampbellLaMotte1983first, Treede1995heat} & measured \\
5 & The gate: an output neuron of the spinal cord receives pain and touch, and an inhibitory interneuron is excited by touch (Fig.~\ref{fig:gate}) & The gate scheme was proposed as a theory. Mice, genetic labeling and ablation of neuron groups: excitatory somatostatin neurons are needed for mechanical pain; inhibitory dynorphin neurons keep input from touch sensors from reaching them, and without these neurons a light touch causes pain. & \cite{MelzackWall1965, Duan2014} & measured \\
6 & Touch dampens pain & Humans, laser pain pulses on the arm: simultaneous touch lowers the pain rating and the ability to discriminate pulse energy; the effect weakens with the distance between the touch point and the pain point within one skin segment. & \cite{Mancini2014touch} & measured \\
7 & Assumption of gate theory: strong pain itself switches off the inhibitory interneuron, and the gate swings open & Marked in the chapter as an assumption of the original theory. The mouse work describes how the gate keeps touch out of the pain channel; switching off of the interneuron by pain is not shown in it. & \cite{MelzackWall1965} & hypothesis \\
8 & The gate~\eqref{eq:gate}: threshold $0.18$ without touch, $0.86$ with touch, $0.11$ at $g=2$; at $\nu=1$ touch lowers the output from $1$ to $0.25$, at $\nu=2$ it has no effect; touch alone does not get through (Fig.~\ref{fig:gatecurves}) & A calculation with \texttt{models/\allowbreak pain\_\allowbreak gate.py} using the weights of the text reproduces all these numbers. & \texttt{models/\allowbreak pain\_\allowbreak gate.py} & model \\
9 & Descending pathways from the brainstem change the gain of the gate in both directions & Rat, medulla, recordings during tail heating: some neurons fire before withdrawal (``on''), others fall silent (``off''); weak stimulation of this area suppresses the reflex. Review: the same circuits both facilitate and inhibit pain transmission, and opioids act through them. & \cite{Fields1983rvm, Fields2004} & measured \\
10 & Expectation of relief dampens pain through the opioid channel & People with acute pain: in those whose pain was reduced by the expectation of relief, blocking the opioid receptors brought the pain back to the level of the others; in the others the blockade did not change the pain. & \cite{Levine1978} & measured \\
11 & The brain chooses the volume of the alarm by the benefit of interrupting & There is no experiment in which the rule for choosing the volume has been measured. & --- & hypothesis \\
12 & Sensitization: after injury the gate output grows and the threshold falls, partly because of the spinal cord; it persists after the damaging stimulus has stopped & Rat: after skin injury the threshold of the flexion reflex falls and the response grows; electrophysiological analysis shows that part of the growth arises in the spinal cord itself. A damaging stimulus causes long-lasting changes in sensitivity that outlast the stimulus. The chapter gives no exact decay time. & \cite{Woolf1983} & measured \\
13 & Sensitization~\eqref{eq:sensitization} is positive feedback; with a strong loop the alarm can latch after healing & Follows from the model~\eqref{eq:sensitization} (an exercise at the end of the chapter). There is no experiment showing that persistent pain after healing is a latched loop of exactly this kind. & derived in the chapter & hypothesis \\
14 & Pain is a negative reward, and learning from it follows the temporal-difference error & Scanner, humans, chains of pain predictors: activity of the ventral striatum and the anterior insula matches temporal-difference learning signals. Monkey: some \nt{DOPA} neurons are excited by an aversive air puff and its predictor, others are inhibited. & \cite{Seymour2004, Matsumoto2009} & measured \\
15 & Pain interrupts the current work & Humans, an electrical pain stimulus and auditory probes: the response to a probe $100\ms$ after pain onset is markedly slowed, and after $1.5\,\text{s}$ there is no longer a difference from control. A review combines such experiments into a model of attentional capture. & \cite{Crombez1996disrupt, EcclestonCrombez1999} & measured \\
16 & Withdrawal is wired in the spinal cord & Rat: neurons of the deep layers of the dorsal horn reproduce the spatial pattern of the withdrawal reflex of individual muscles; they cannot be excited antidromically from the cervical cord, so they are spinal interneurons. & \cite{Schouenborg1995wr} & measured \\
\bottomrule
\end{longtable}}
